\documentclass[12pt]{article}
\PassOptionsToPackage{table}{xcolor}
\usepackage[utf8]{inputenc}
\usepackage[T1]{fontenc}
\usepackage[italian]{babel}
\usepackage{multicol}
\usepackage{enumitem}
\usepackage{array}
\usepackage{xcolor}
\usepackage{graphicx}
\usepackage{tikz}
\usetikzlibrary{
    shapes.geometric,
    arrows.meta,
    positioning,
    calc,
    patterns,
    decorations.pathmorphing
}
\usepackage[margin=2cm]{geometry}
\usepackage{amsmath,amssymb}
\setlength{\parindent}{0pt}
\definecolor{headercolor}{RGB}{220,220,220}
\newcounter{quesito}
\setcounter{quesito}{0}
\newcommand{\nuovoquesito}{\stepcounter{quesito}\textbf{Domanda \arabic{quesito}.}}
\newcommand{\itemdomanda}[2]{%
    \par\noindent
    \begin{minipage}[t]{\linewidth}
        \nuovoquesito \\ #1
        \begin{enumerate}[label=(\Alph*), nosep, leftmargin=*]
            #2
        \end{enumerate}
    \end{minipage}
    \vspace{10pt}
}
\begin{document}
\section*{Simulazione}
\hfill Data: 16 apr 2026

\itemdomanda{Se dopodomani sarà giovedì, che giorno era ieri?}{
    \item Lunedì
    \item Martedì
    \item Domenica
    \item Mercoledì
    \item Sabato
}

\itemdomanda{Date le premesse:
    \begin{enumerate}
        \item Nessun cane terrier gironzola tra i segni dello zodiaco.
        \item Tra quelli che non gironzolano tra i segni dello zodiaco, nessuno è una cometa.
        \item Solo i cani terrier hanno la coda arricciata.
    \end{enumerate}
    Si può concludere che:}{
    \item Nessuna cometa ha la coda arricciata
    \item Qualche cometa ha la coda arricciata
    \item Tutte le comete sono cani terrier
    \item Chi ha la coda arricciata è una cometa
    \item I cani terrier sono comete
}

\itemdomanda{Volendo disporre i numeri $28, 29, 36, 43, 55$ in modo che i dispari occupino una posizione dispari ed i pari occupino una posizione pari, in quanti modi diversi si può operare?}{
    \item 3
    \item 24
    \item 12
    \item 5
    \item 6
}

\itemdomanda{Nonno Peperino non ricorda più la combinazione del suo forziere elettronico. Egli ricorda solo che: è di quattro cifre distinte fra 0 e 9; non vi compare il 4; la terza cifra è la metà della quarta; le cifre sono in ordine crescente dalla prima all'ultima. Qual è il minimo numero di tentativi che Nonno Peperino deve fare per essere sicuro di aprire il forziere?}{
    \item 3
    \item 4
    \item 5
    \item 6
    \item 2
}

\itemdomanda{Due giocatori, Aldo e Bruno, a turno dispongono su una scacchiera $3 \times 3$, una per volta, pedine identiche tutte nere. Vince il primo giocatore che riesce a completare un terzetto in una fila orizzontale, verticale o una delle due diagonali. Quante sono le mosse con cui può iniziare la partita il primo giocatore (Aldo) in modo da potersi garantire la vittoria indipendentemente da come giocherà Bruno?}{
    \item 8
    \item 1
    \item 0, nel senso che Bruno può sempre rispondere in modo opportuno e garantirsi la vittoria
    \item 5
    \item 9
}

\itemdomanda{Due giocatori prendono a turno dei sassolini con l'unica regola che non se ne possono prendere né 4 né 8. Vince quel giocatore che riesce a prendere l'ultimo sassolino. Se inizialmente i sassolini sono 8, quanti ne deve prendere il primo giocatore per potersi garantire la vittoria, supponendo che nelle mosse successive ogni giocatore non commetta errori?}{
    \item Qualunque numero prenda, vincerà sempre
    \item Qualunque numero prenda, perderà sempre
    \item 1
    \item 2
    \item 3
}

\itemdomanda{Quale fra le seguenti affermazioni è falsa?}{
    \item Chi canta è allegro. Andrea non canta, dunque Andrea è triste.
    \item Nessun londinese è italiano; tutti i londinesi parlano inglese, ma non è vero che nessun italiano parla inglese.
    \item Un quadrato è sempre un trapezio.
    \item Ogni vigile ha un fischietto. Carlo non ha un fischietto, dunque Carlo non è un vigile.
    \item "Chi canta è allegro. Andrea non canta, dunque Andrea è triste." non è vero.
}

\itemdomanda{Premesso che:
    \begin{itemize}
        \item chi ascolta musica rock o blues non è stonato
        \item Agenore non è stonato
        \item chi ascolta blues non vince al Lotto
    \end{itemize}
    quale tra le seguenti conclusioni \textbf{NON} si può trarre dalle precedenti premesse?}{
    \item È impossibile che Agenore ascolti blues
    \item Uno stonato non ascolta rock
    \item È possibile che Agenore non vinca al Lotto
    \item Chi vince al Lotto non ascolta blues
    \item Non è escluso che Agenore ascolti rock
}

\itemdomanda{Dalla frase "O mangi questa minestra o ti salti dalla finestra" (interpretata come O esclusivo), si deduce che:}{
    \item Se mangi la minestra, puoi anche saltare dalla finestra
    \item Se non salti dalla finestra, allora devi mangiare la minestra
    \item Puoi fare entrambe le cose contemporaneamente
    \item Non puoi evitare entrambe le situazioni
    \item Se non mangi la minestra, non è detto che tu debba saltare
}

\itemdomanda{Data la frase "Se piove, porto l'ombrello", quale delle seguenti è certamente vera?}{
    \item Se porto l'ombrello, allora piove
    \item Se non porto l'ombrello, allora non piove
    \item Se non piove, non porto l'ombrello
    \item Piove solo se porto l'ombrello
    \item Porto l'ombrello se e solo se piove
}

\itemdomanda{Quale fra le seguenti affermazioni è sicuramente \textbf{falsa}?}{
    \item Chi respira è vivo. Piero non respira, dunque Piero è morto
    \item Nessun parigino è italiano; tutti i parigini parlano francese; ma non è vero che nessun italiano parla francese.
    \item Un quadrato è sempre un rombo
    \item Ciò che è scritto in A. è falso
    \item Ogni professore ha un registro. Mario non ha registro, dunque Mario non è professore
}

\itemdomanda{In quale di queste sequenze letterarie viene mantenuta l'alternanza vocale-consonante?}{
    \item ODERAXETUIXIDIC
    \item ODERAXETUPAFESTXIDIC
    \item ODERAXETZIXIDIC
    \item ODERAXETUPAFESIXIDIC
    \item ODORAXITUPAFASTXIDIC
}

\end{document}
